\documentclass{article}
\usepackage[utf8]{inputenc}
\usepackage{array}
\newcolumntype{C}[1]{>{\centering\arraybackslash}m{#1}}
\usepackage{amsmath}
\usepackage{relsize}
\usepackage{graphicx}
\usepackage{listings}
\usepackage{algorithm}
\usepackage{algpseudocode}
\usepackage{hyperref}
\usepackage{tikz}
\usepackage{tocloft}
\usepackage[explicit]{titlesec}
\usepackage[margin=1in]{geometry}
\usetikzlibrary{trees}

\renewcommand{\contentsname}{Indice}
\renewcommand\cftaftertoctitle{\hfill\mbox{}}

\titleformat{\section}{\normalfont\Large\bfseries}{}{0em}{#1\ \thesection}

\lstset{
    language=C++,
    basicstyle=\ttfamily,
    keywordstyle=\color{blue}\ttfamily,
    stringstyle=\color{red}\ttfamily,
    commentstyle=\color{green}\ttfamily,
    morecomment=[l][\color{magenta}]{\#},
    numbers=left,
    stepnumber=1,
    showstringspaces=false,
    tabsize=1,
    breaklines=true,
    breakatwhitespace=false,
}

\usepackage[noabbrev]{cleveref}

\crefformat{section}{\S#2#1#3}
\crefformat{subsection}{\S#2#1#3}
\crefformat{subsubsection}{\S#2#1#3}
\crefformat{figure}{#2Figura~#1#3}
\crefformat{footnote}{#2\footnotemark[#1]#3}
\crefformat{table}{#2Tabella~#1#3}
\crefformat{equation}{#2Equazione~#1#3}
\crefformat{chapter}{#2Capitolo~#1#3}

\begin{document}

\begin{titlepage}
    \begin{center}
        \vspace*{1cm}

        \Huge
        \textbf{Corso di Algoritmi e Strutture Dati}

        \vspace{0.5cm}
        \LARGE
        Homework

        \vfill
        \vspace{0.8cm}

        \Large
        Capasso Giuseppe - M63001498

    \end{center}
\end{titlepage}

\tableofcontents
\newpage

\addcontentsline{toc}{section}{Richiami teorici}
\subsection*{Richiami teorici}
Nello svolgimento degli esercizi sono stati utilizzati i seguenti concetti teorici.

\paragraph{Notazione $O$}
Data una funzione $g(n)$, $O(g(n))$ indica un insieme di funzioni $f(n)$ in cui:
\begin{equation}
    O(g(n)) = \{f(n): \exists c > 0, \exists m > 0: f(n) \leq c\cdot g(n), \; \forall n \geq m\}
\end{equation}

\paragraph{Notazione $\Omega$}
Data una funzione $g(n)$, $\Omega(g(n))$ indica un insieme di funzioni $f(n)$ in cui:
\begin{equation}
    \Omega(g(n)) = \{f(n): \exists c > 0, \exists m > 0: f(n) \geq c\cdot g(n), \; \forall n \geq m\}
\end{equation}

\paragraph*{Notazione $\Theta$}
Data una funzione $g(n)$, $\Theta(g(n))$ indica un insieme di funzioni $f(n)$ in cui:
\begin{equation}
    \Theta(g(n)) = \{f(n): \exists c_1, c_2 > 0, \exists m > 0: c_1\cdot g(n) \leq f(n) \leq c_2\cdot g(n), \; \forall n \geq m\}
    \label{eq:theta}
\end{equation}

\paragraph{Relazioni di dominanza elementari}
Dette $p < r$, $h < k$, $a < b$, valgono le seguenti relazioni per i problemi computazionali:
\begin{align*}
     O(1) &\subset O(log^r(n)) \subset O(log^p(n)) \subset O(n^h) \subset O(n^h \cdot log^r(n))\\ 
     &\subset O(n^h \cdot log^p(n)) \subset O(n^k) \subset O(a^n) \subset O(b^n) \subset O(n!)
\end{align*}

\paragraph{Proprietà di $O$}
In termini di $O$, valgono le seguenti proprietà:
\begin{itemize}
    \item \textbf{Sequenza di algoritmi}
    \begin{equation}
        f(n) + g(n) = O(max(f(n), g(n)))
        \label{eq:alg-seq}
    \end{equation}
    \item \textbf{Eliminazione delle costanti}
    \begin{equation}
        O(c f(n)) = O(f(n))
        \label{eq:const-elim}
    \end{equation}
    \item \textbf{Cicli annidati}
    \begin{equation}
        O(f(n)) \cdot O(g(n)) = O(f(n) \cdot g(n))
        \label{eq:nested-loops}
    \end{equation}

    \item \textbf{Riflessività}
    \begin{equation}
     O(f(n)) = f(n)
     \label{eq:reflex}
    \end{equation}
\end{itemize}

\paragraph{Serie matematiche note}
\begin{equation}
    \mathlarger{\sum}_{i=0}^{\infty}(x^i) = \frac{1}{1-x} \quad con\; 0\leq x<1
    \label{eq:geometric_sum}
\end{equation}

\begin{equation}
    \mathlarger{\sum}_{i=0}^{\infty}(i\cdot x^i) = \frac{x}{(1-x)^2} \quad con\; 0\leq x<1
    \label{eq:sum}
\end{equation}

\paragraph{Teorema dell'esperto}\label{s:master-th}
Data una ricorrenza nella forma $T(n) = aT\left(\frac{n}{b}\right) + f(n))$ e, detto $\alpha = \log_b(a)$, il teorema dell'esperto permette di calcolare facilmente un limite alla complessità delle ricorrenze, secondo la tabella sottostante:

\begin{center}
    \begin{tabular}{| C{0.75cm} | C{5.0cm} | C{5.0cm} |}
        \hline
        Caso & Condizione &  Risultato \\ [2ex]
        \hline\hline
        1 & $f(n) = O(n^{\alpha-\epsilon}), \epsilon > 0$ & $T(n) = \Theta(n^\alpha)$ \\  [2ex]
        \hline
        2 & $f(n) = \Theta(n^\alpha)$ & $T(n) = \Theta(n^\alpha\cdot\log(n))$ \\ [2ex]
        \hline
        3 & $f(n) = \Omega(n^{c+\epsilon})$ e $a f(n/b)\leq cf(n)\; con\; c<1$ & $T(n) = \Theta(f(n))$ \\ [2ex]
        \hline
    \end{tabular}
\end{center}

\paragraph{Ricorrenze lineari di ordine costante}\label{s:linaer-rec}
Siano $a_1, a_2, \dots, a_h$ costanti intere non negative con $h > 0$, $c, b \in R$ tali che $c > 0$ e $\beta \geq 0$ e sia definita la ricorrenza $T(n)$:
\begin{equation*}
    T(n) = \begin{cases}
        \sum_{1\leq i \leq h} (a_i T(n-i) + cn^\beta & n \geq m\\
        \Theta(1) & n < m
    \end{cases}
\end{equation*}

Detto $a = \sum_{1\leq i \leq h}(a_i)$, allora è possibile affermare che:
\begin{equation*}
    T(n) = \begin{cases}
        \Theta(n^{\beta+1}) &\text{se} \; a = 1\\
        \Theta(a^n n^{\beta}) &\text{se}\; a\geq 1\\
    \end{cases}
\end{equation*}
\clearpage
\section{Homework}
\subsection{Notatione asintotica}

Per ognuna delle seguenti affermazioni, si dica se essa è sempre vera, mai vera, o a volte
vera, per funzioni asintoticamente non-negative. Se la si considera sempre vera o mai vera,
si spieghi il perché. Se è a volte vera, si dia un esempio per cui è vera e uno per cui è falsa.

\subsection*{A. $f(n) = O(f(n)^2)$}
La relazione risulta sempre vera in accordo con l'ipotesi di funzioni asintoticamente non negative. Infatti, applicando la definizione \textit{Big O}, si ottiene:
\begin{equation*}
    \exists c > 0, \exists m \geq 0: \; f(n)\leq c f^2(n)\; \forall n\geq m
\end{equation*}
che risulta sempre vera per $c=1$ per tutte le funzioni monotone crescenti.

\subsection*{B. $f(n) + O(f(n)) = \Theta(f(n))$}
Applicando la proprietà di riflessività secondo la quale $O(f(n)) = f(n)$, è possibile scrivere:
\begin{equation*}
    f(n) + O(f(n)) = 2f(n) = \Theta(f(n))
\end{equation*}
A questo punto, applicando nuovamente la proprietà di riflessività e quella per l'eliminazione delle costanti , si può scrivere che:
\begin{equation*}
   2f(n) = \Theta(2f(n)) = \Theta(f(n))
\end{equation*}
Quindi l'affermazione risulta vera nel caso in cui $f(n)$ sia una funzione asintoticamente non-negativa.

\subsection*{C. $f(n) = \Omega(g(n))$ e $f(n) = o(g(n))$}
L'affermazione risulta falsa perché il simbolo \textit{o-piccolo} non include l'uguaglianza. Infatti, applicando la definizione \textit{Omega}, si ottiene:
\begin{equation*}
    \exists c > 0, \exists m \geq 0: \; f(n)\geq c\cdot g(n)\; \forall n\geq m
\end{equation*}
Mentre, usando quella di \textit{o-piccolo}, si ha che:
\begin{equation*}
    \exists c > 0, \exists m \geq 0: \; f(n) < c \cdot g(n)\; \forall n\geq m
\end{equation*}
La condizione \textit{and} non è mai verificata per nessun valore di $c$ in quanto o $f(n) \geq g(n)$ o $f(n) < g(n)$.
Invece, se al posto di $o(g(n))$ ci fosse stato $O(g(n))$, allora la relazione sarebbe stata vera in quanto se $f(n) = \Omega (g(n))$ e $f(n) = O(g(n))$ allora $f(n) = \Theta(g(n))$.

\clearpage
\subsection{Notazione asintotica e crescita delle funzioni}
Per ognuna delle seguenti coppie di funzioni $f(n)$ e $g(n)$, trovare una appropriata costante positiva $c$ tale che $f(n) \leq c \cdot g(n)$ per tutti i valori di $n > 1$.

\subsection*{A. $f(n) = n^2 + n + 1$, $g(n) = 2n^3$}
Sfruttando il fatto che $n > 1$, è possibile scrivere che:
\begin{align*}
    f(n) &= n^2 + n + 1\\
  &\leq n^3 + n + 1\\
  &\leq n^3 + n^3 + 1\\
  &\leq n^3 + n^3 + n^3\\
  &= 3n^3\\
  &\leq c\cdot g(n) \\
  &=c \cdot 2n^3
\end{align*}

Quindi, bisogna risolvere la relazione $3n^3 \leq c \cdot 2n^3$ che risulta vera per $c \geq \frac{3}{2}$.

\subsection*{B. $f(n) = n\sqrt{n} + n^2$, $g(n) = n^2$}
Essendo $n > 1$, è possibile scrivere che:
\begin{align*}
    f(n) &= n\sqrt{n} + n^2\\
  &\leq n^2 + n^2\\
  &= 2n^2\\
  &\leq c\cdot g(n) \\
  &=c \cdot n^2
\end{align*}

Da cui si ottiene la disequazione $2n^2 \geq cn^2$ che risulta vera per $c \geq 2$.

\subsection*{C. $f(n) = n^2 - n + 1$, $g(n) = \frac{n^2}{2}$}
Essendo $n > 1$, è possibile maggiorare $f(n)$ come segue:
\begin{align*}
    f(n) &= n^2 - n + 1\\
  &\leq n^2 +1 \quad \text{essendo $n > 1$, basta eliminare $-n$ per rendere vera la maggiorazione}\\
  &\leq n^2 + n^2 \\
  &= 2n^2\\
  &\leq c\cdot g(n) \\
  &=c \cdot \frac{n^2}{2}
\end{align*}

Quindi il problema si riduce a risolvere la disequazione $2n^2 \leq c\cdot \frac{n^2}{2}$ che risulta vera per $c \geq 4 $.

\clearpage
\subsection{Complessità}
Dimostrare che per qualsiasi costante reale $a$ e $b$, con $b > 0$, $(n+a)^b = \Theta(n^b)$.\\
\paragraph{Dimostrazione}
Per dimostrare la relazione, si osserva che l'operazione\\ $(n+a)^b$ genera sempre un polinomio di grado $b$ che può essere scritto nella forma:
\begin{equation}
    f(n) = n^b + c_{b-1}n^{b-1}+\ldots+c_1n+c_0
    \label{eq:polinomio}
\end{equation}

Quindi dimostrare che $(n+a)^b = \Theta(n^b)\;,b > 0$ equivale a provare che $f(n)= \Theta(n^b)$, dove con $f(n)$ si indica il polinomio in \cref{eq:polinomio}.

Applicando la definizione di $\Theta$ (\cref{eq:theta}), si procede con il dimostrare prima il limite superiore e poi quello inferiore.

\paragraph{Limite superiore} Bisogna dimostrare che $\exists c_2 > 0, \exists m > 0: f(n) \leq c_2\cdot n^b, \forall n \geq m$. Supponendo $n \geq 1$, si può scrivere che:
\begin{align*}
    f(n)&= n^b+c_{b-1}n^{b-1}+\ldots+c_1n+c_0\\
    &\leq n^b+|c_{b-1}|n^{b-1}+\ldots+|c_1|n+|c_0|\\
    &\leq (1+|c_{b-1}|+\ldots+|c_1|+|c_0|)n^b\\
    &\leq c_2\cdot n^b
\end{align*}

La relazione risulta vera per $c_2 \geq 1+|c_{b-1}|+\ldots+|c_1|+|c_0| $

\paragraph{Limite inferiore} Analogamente al limite superiore, bisogna dimostrare che $\exists c_2 > 0, \exists m > 0: f(n) \geq c_1\cdot n^b, \forall n \geq m$, si può scrivere che:

\begin{align*}
    f(n)&= n^b+c_{b-1}n^{b-1}+\ldots+c_1n+c_0\\
    &\geq n^b-|c_{b-1}|n^{b-1}-\ldots-|c_1|n-|c_0|\\
    &\geq (n-|c_{b-1}|-\ldots-|c_1|-|c_0|)n^{(b-1)}\\
    &\geq c_1\cdot n^b
\end{align*}

La disequazione è soddisfatta se $c_1 \leq (1 - \frac{|c_{b-1}|+\ldots+|c_1|+|c_0|}{n}) $

Quindi avendo dimostrato che esistono $c_1,c_2 > 0$ tale che $c_1\cdot n^b\leq f(n) \leq c_2 \cdot n^b$, allora è possibile affermare che $f(n) = \Theta(n^b)$.

\clearpage
\subsection{Ricorrenze}
Fornire il limite inferiore e superiore per $T(n)$ nella seguente ricorrenza, usando il metodo dell’albero delle ricorrenze ed il teorema dell’esperto se applicabile. Si fornisca il limite più stretto possibile giustificando la risposta.

\subsection*{A. $T(n) = 2T(\frac{n}{3})+n\cdot \log n$}
\paragraph{Soluzione con albero di ricorrenza}\mbox{}\\

Di seguito, è mostrato l'albero di ricorrenza generato da $T(n)$.

\begin{tikzpicture}[level 1/.style={sibling distance=5.6cm},level 2/.style={sibling distance=3cm},level 3/.style={sibling distance=2cm}]
	\node {$cn$}[edge from parent fork down]
		child { node {$c(n/3)$}
			child { node {$c(n/9)$}
			    child {node{...}}
			    child {node{...}}
			    }
			child { node {$c(n/9)$}
			    child {node{...}}
			    child {node{...}}
			    }
			}
		child { node {$c(n/3)$}
			child { node {$c(n/9)$}
			    child {node{...}}
			    child {node{...}}
			    }
			child { node {$c(n/9)$}
			    child {node{...}}
			    child {node{...}}
			    }
			}
	;
\end{tikzpicture}\mbox{}\\
L'albero di ricorrenza può essere schematizzato secondo la tabella sottostante, che rappresenta per ogni livello qual è la dimensione del problema, il costo di ciascuna chiamata, il numero di chiamate ricorsive e il costo totale di livello (ottenuto moltiplicando il numero di chiamate per il costo di chiamate).

Al livello \textit{i-esimo} ci sono $2^i$ chiamate (nodi nell'albero) ciascuno dei quali con costo $n/3^i$. La ricorsione termina quando si raggiunge il caso base $T(1)$ ovvero quando $l = n/3^i = 1 \rightarrow i = \log_3 n$.

\begin{center}
    \begin{tabular}{| C{1.5cm} | C{1.7cm} | C{2.7cm} | C{1.5cm} | C{3.6cm} |}
        \hline
        Livello & Dimensione &  Costo Chiamata & Num. Chiamate & Costo Livello \\
        \hline\hline
        $0$ & $n$ & $n \log n$ & $1$ & $n \log n$\\
        \hline
        $1$ & $(n/3)$ & $n/3 \log (n/3)$ & $2$ & $2 (n/3) \log (n/3)$\\[7pt]
        \hline
        $2$ & $(n/3^2)$ & $n/3^2 \log (n/3^2)$ & $2^2$ & $2^2 (n/3^2) \log (n/3^2)$\\[7pt]
        \hline
        \dots & \dots & \dots & \dots & \dots \\
        \hline
        $i$ & $(n/3^i)$ & $n/3^i \log (n/3^i)$ & $2^i$ & $2^i (n/3^i) \log (n/3^i)$\\[7pt]
        \hline
        \dots & \dots & \dots & \dots & \dots \\
        \hline
        $l-1$ & $(n/3^{l-1})$ & $n/3^{l-1} \log (n/3^{l-1})$ & $2^{l-1}$ & $2^{l-1}(n/3^{l-1}) \log (n/3^{l-1})$\\[7pt]
        \hline
        $l = \log_3 n$ & $1$ & $T(1)$ & $2^l$ & $2^l$\\[7pt]
        \hline
    \end{tabular}
\end{center}

Per ottenere il costo totale dell'algoritmo bisogna sommare i costi di tutti i livelli:
\begin{equation*}
    T(n) = \sum_{i=0}^{\log_3n-1} \left(2^i \frac{n}{3^i} \log \frac{n}{3^i}\right) + 2^{\log_3 n}
\end{equation*}
Applicando la proprietà $\log(a/b) = \log(a) - \log (b)$, si ottiene:

\begin{equation*}
    T(n) = n \left[\log(n)\sum_{i=0}^{\log_3n-1} \left(\frac{2}{3}\right)^i - \log3 \sum_{i=0}^{\log_3n-1} i \cdot \left(\frac{2}{3}\right)^i \right] + 2^{\log_3 n}
\end{equation*}

Applicando le sommatorie note (\cref{eq:geometric_sum} e \cref{eq:sum}) e la proprietà $a^{\log_b c} = c ^{\log_b a}$, si ottiene:

\begin{align*}
     T(n) &= n \left[\log(n)\sum_{i=0}^{\log_3n-1} \left(\frac{2}{3}\right)^i - \log3 \sum_{i=0}^{\log_3n-1} i \cdot \left(\frac{2}{3}\right)^i \right] + 2^{\log_3 n}\\
    &\leq n \log(n) \left(\frac{1}{1-2/3}\right) - n \frac{2/3}{(1-2/3)^2} + n^{\log_3 2} \\
    &= 3n\log(n) - 6n + n^{\log_3 2}\\
\end{align*}

Usando la notazione $big-O$, applicando la proprietà di eliminazione delle costanti e le relazioni di dominanza elementari, si ottiene;

\begin{align*}
    T(n) &= 3n\log(n) - 6n + n^{\log_3 2}\\
    &=O(3n\log n  - 6n + n^{\log_3 2})\\
\end{align*}

Quindi il costo dell'algoritmo è $T(n) = O(n\log n)$.

\paragraph{Soluzione con il teorema dell'esperto}
Secondo il teorema dell'esperto, si individuano i valori $a = 2$ e $b = 3$ da cui si ottiene che $\alpha = \log_b a = \log_3 2 \approx 0.63$; mentre $f(n) = n \log n$. Possono essere esclusi i casi 1 e 2 del teorema dell'esperto poiché $f(n) = n \log n \neq O(n^{0.63-\epsilon})$ (caso 1) in quanto $\epsilon > 0$ e sicuramente \textbf{non} vale la relazione  $f(n) = n \log n = \Theta(n^{0.63})$.
Invece, per il caso 3, si può provare che  $$f(n) = n \log n = \Omega(n^{0.63+\epsilon})$$ che risulta vera per $\epsilon < 1 - \alpha \approx 0.37$. Inoltre, bisogna provare che:
\begin{equation*}
    \exists c \leq 1: \exists m \geq 0: a f\left(\frac{n}{b}\right) \leq c\cdot f(n)
\end{equation*}

\begin{align*}
    af\left(\frac{n}{b}\right) &= 2 \frac{n}{3} \log \left(\frac{n}{3}\right)\\
  &=  2 \frac{n}{3} \log n - \log 3\\
  &\leq 2 \frac{n}{3} \log n\\
  &\leq c \cdot n  \log n
\end{align*}

Da cui si ricava che $c \geq \frac{2}{3}$. Quindi, dimostrate le ipotesi del caso 3 del teorema dell'esperto, si può affermare che $T(n) = \Theta(n \cdot \log n)$.

\subsection*{B. $T(n) = 3T(\frac{n}{5})+\log^2 n$}
\paragraph{Soluzione con albero di ricorrenza}\mbox{}\\

Di seguito, è mostrato l'albero di ricorrenza generato da $T(n)$ (semplificato per motivi di spazio).

\begin{tikzpicture}[level 1/.style={sibling distance=3.5cm},level 2/.style={sibling distance=2.5cm},level 3/.style={sibling distance=0.85cm}]
	\node {$cn$}[edge from parent fork down]
        
    child { node {$c(n/5)$}
        child { node {$c(n/25)$} 
              child {node{...}}
			    child {node{...}}
              child {node{...}}
       }
        child { node {$c(n/25)$}
          child {node{...}}
			child {node{...}}
          child {node{...}}
        }
        child { node {$c(n/25)$}
          child {node{...}}
			child {node{...}}
          child {node{...}}
        }
    }
    child { node {$c(n/5)$}}
    child { node {$c(n/5)$}};
\end{tikzpicture}\\
L'albero di ricorrenza è schematizzato nella tabella sottostante.
\begin{center}
    \begin{tabular}{| C{1.5cm} | C{1.7cm} | C{2.7cm} | C{1.5cm} | C{3.6cm} |}
        \hline
        Livello & Dimensione &  Costo Chiamata & Num. Chiamate & Costo Livello \\ 
        \hline\hline
        $0$ & $n$ & $n \log^2 n$ & $1$ & $\log^2 n$\\
        \hline
        $1$ & $(n/5)$ & $\log^2 (n/5)$ & $3$ & $3 \log (n/5)$\\[7pt]
        \hline
        $2$ & $(n/5^2)$ & $\log^2 (n/5^2)$ & $3^2$ & $3^2 \log^2 (n/5^2)$\\[7pt]
        \hline
        \dots & \dots & \dots & \dots & \dots \\
        \hline
        $i$ & $(n/5^i)$ & $\log^2 (n/5^i)$ & $3^i$ & $3^i \log^2 (n/5^i)$\\[7pt]
        \hline
        \dots & \dots & \dots & \dots & \dots \\
        \hline
        $l-1$ & $(n/5^{l-1})$ & $\log^2 (n/5^{l-1})$ & $3^{l-1}$ & $3^{l-1}\log^2 (n/5^{l-1})$\\[7pt]
        \hline
        $l = \log_5 n$ & $1$ & $T(1)$ & $3^l$ & $3^l$\\[7pt]
        \hline
    \end{tabular}
\end{center}
Al livello \textit{i-esimo} ci sono $3^i$ nodi ciascuno dei quali con costo $n/5^i$. Il caso base viene raggiunto quando  $n/5^i = 1 \rightarrow i = \log_5 n$.  Il costo di ciascuna chiamata è dato da $\log^2 n$ e, quindi, il costo al livello \textit{i-esimo} è $3^i\log^2 n/5^i$. Quindi il costo totale dell'algoritmo è dato da:
\begin{equation*}
    \sum_{i=0}^{\log_5 n-1} \left(3^i\log^2 (n/5^i)\right) + 3^{\log_5 n}
\end{equation*}

Attraverso passaggi algebrici analoghi a quelli dell'esercizio precedente si ottiene:
\begin{align*}
    T(n) &= \sum_{i=0}^{\log_5 n-1} (3^i\log^2 n) -  \sum_{i=0}^{\log_5 n-1} (3^i\log^2 (5^i)) + n^{\log_5 3} \\
    &= \log^2 n \sum_{i=0}^{\log_5 n-1} (3^i) - \log^2 5 \sum_{i=0}^{\log_5 n-1} (i 3^i) + n^{\log_5 3}
\end{align*}

È possibile maggiorare $T(n)$ ottenuta eliminando il termine negativo centrale:

\begin{align*}
    T(n) &= \log^2 n \sum_{i=0}^{\log_5 n-1} (3^i) - \log^2 5 \sum_{i=0}^{\log_5 n-1} (i 3^i) + n^{\log_5 3}\\
    &\leq \log^2 n \sum_{i=0}^{\log_5 n-1} (3^i) + n^{\log_5 3}
\end{align*}

Essendo $\log_5 n$ una quantità finita, si può affermare che la sommatoria al primo termine ha un valore finito $M$. Sapendo che $\log_5 3 \approx 0.68$ si può scrivere che:

\begin{align*}
    T(n) &\leq \log^2 n \sum_{i=0}^{\log_5 n-1} (3^i) + n^{\log_5 3}\\
    &= M\log^2 n + n^{0.68}\\
    &= O(M\log^2 n + n^{0.68})
\end{align*}

Applicando la proprietà \cref{eq:alg-seq} e la relazione di dominanza si può affermare che:
\begin{equation*}
    T(n) = O(n^{0.68})
\end{equation*}


\paragraph{Soluzione con il teorema dell'esperto}
Secondo il teorema dell'esperto, si individuano i valori $a = 3$ e $b = 5$ da cui si ottiene che $\alpha = \log_b a = \log_5 3 \approx 0.68$; mentre $f(n) = \log^2 n$. In questo caso, è possibile applicare il caso 1 del teorema dell'esperto, in quanto $f(n) = O(n^{(0.68-\epsilon)}), \epsilon > 0$. Infatti, bisogna verificare che $$f(n) = \log^2 n \leq c n^{0.68}$$ come mostrato nei seguenti passaggi:
\begin{align*}
    \log^2 n &\leq c n^{0.68} \quad \text{si elevano entrambi i membri alla $\frac{1}{2}$}\\
    \log n &\leq c_1 n^{0.34} \quad \text{con $c_1 =\sqrt{c} $}\\
    2^{\log n} &\leq 2^{c_1 n^{0.34}} \quad \text{si elimina il logaritmo}\\
    n &\leq n^{c_1 2^{0.34}} \\
    n &\leq n^{c_11.266}
\end{align*}

che risulta vera per $c_1 \geq 1$. Quindi applicando il caso 1 del teorema dell'esperto, si ottiene che $T(n) = \Theta(n^{\log_5 3}) \approx \Theta(n^{0.68})$.

\clearpage

\subsection{Problema 1}
Si implementi un algoritmo di ordinamento che sfrutta l’inserimento e la visita in un albero
binario di ricerca. Dato un vettore di n numeri interi in input, l’algoritmo procede prima ad inserire i numeri in un albero binario di ricerca (usando ripetutamente TREE-INSERT per
inserire i numeri uno alla volta), e poi stampa i numeri in ordine con un attraversamento in
ordine simmetrico dell’albero. Si analizzi la complessità nel caso peggiore e nel caso
migliore di per questo algoritmo di ordinamento.

\paragraph{Soluzione}
L'implementazione prevede la creazione di una struttura dati \textit{Node} che rappresenta il nodo di un albero al cui interno sono memorizzati i riferimenti al figlio destro, quello sinistro e la chiave, come mostrato nel listato di seguito.

\begin{lstlisting}[caption={Nodo di un albero binario di ricerca},captionpos=b]
struct node {
    node *left;
    node *right;
    int key;
};

\end{lstlisting}
Sono state progettate 2 funzioni \textit{insert} e \textit{in\_order\_print} per inserire gli elementi nell'albero e stamparli in ordine da sinistra verso destra.

\begin{lstlisting}[caption={Procedura di inserimento in un BST},captionpos=b]
void insert(struct node **root, struct Node *z) {
  struct node *y = nullptr;
  struct node *x = *root;

  while (x != nullptr) {
    y = x;
    if (z->key < x->key) {
      x = x->left;
    } else {
      x=x->right;
    }
  }

  if (y == nullptr) {
    *root = z;
  } else if (z->key < y->key) {
    y->left = z;
  } else {
    y->right = z;
  }

}

void in_order_print(struct node *root) {
  if (root != nullptr) {
    in_order_print(root->left);
    cout << root->key << ' ';
    in_order_print(root->right);
  }
}

\end{lstlisting}
Infine, l'algoritmo risolutivo (semplificato) è riportato di seguito:

\begin{lstlisting}[caption={Ordinamento di un albero mediante insert},captionpos=b]
void solve(const vector<int>& a) {
  struct node *tree = nullptr;
  for (int i = 0; i < a.size(); ++i) {
    insert(&tree, new Node(a[i]));
  }

  in_order_print(tree);
}

\end{lstlisting}

\paragraph{Complessità}
Detta $n$ la dimensione dell'input e posto $h = \log n$, è possibile calcolare la complessità dell'algoritmo utilizzando le proprietà dei cicli annidati e algoritmi in sequenza(\cref{eq:nested-loops} e \cref{eq:alg-seq}). In particolare, sapendo che l'inserimento all'interno di un \textbf{BST} ha complessità $O(h)$ e la visita in ordine è $\Theta(n)$, si può affermare che il \textit{loop} a riga \textit{3-5} ha una complessità $O(n\log n)$. Quindi il costo totale, nel caso medio, dell'algoritmo è $T(n) = O(n\log n) + \Theta(n) = O(n \log n)$. Nel caso pessimo, bisogna osservare che un albero non bilanciato come il BST può degenerare in una \textit{linked list} (ad esempio quando il vettore di input è già ordinato) e, pertanto, la procedura di inserimento diventa $O(n)$. Quindi, nel caso pessimo la complessità diventa $T(n) = O(n \cdot n) + \Theta(n) = O(n^2)$.
\clearpage

\subsection{Problema 2}
Si implementi un algoritmo che, a partire da un vettore di n numeri interi in input, costruisce un heap chiamando ripetutamente la procedura MAX-HEAP-INSERT per inserire gli elementi nell’heap. L’algoritmo di costruzione ha il seguente
pseudocodice:

\begin{algorithm}
    \begin{algorithmic}
        \Procedure{BUILD-MAX-HEAP\_v2}{$A$}
            \State $A.heap\_size \gets 1$
            \For{$i=2\texttt{ to } A.length$}
                \State $MAX-HEAP-INSERT(A, A[i])$
            \EndFor
        \EndProcedure
    \end{algorithmic}
\end{algorithm}

Si confronti la BUILD-MAX-HEAP vista a lezione con la BUILD-MAX-HEAP\_v2, in
particolare: le due procedure creano sempre lo stesso \textit{heap} se vengono eseguite con lo stesso \textit{array} di input? Dimostrare che lo fanno o fornite un controesempio.
Dimostrare che, nel caso peggiore,\textit{BUILD-MAX-HEAP\_v2} richiedo un tempo $\Theta(n\log n)$ per costruire un \textit{heap} di n elementi.

\paragraph{Soluzione} Per la soluzione si è implementata una struttura dati \textit{heap} che ha la seguente interfaccia:

\begin{lstlisting}[caption={Interfaccia di un max heap},captionpos=b]
struct heap {
   size_t heap_size;
   size_t capacity;
   int *keys;

   void max_heapify(const int i);
   void heap_increase_key(int i, int key);
   void max_heap_insert(const int k);

}
\end{lstlisting}
Per implementare la procedura \textit{BUILD-MAX-HEAP\_v2}, è stato necessario implementare \textit{max\_heap\_insert} e \textit{heap\_increase\_key} riportate di seguito:

\begin{lstlisting}[caption={Implementazione procedure di inserimento in un max heap},captionpos=b]
void heap_increase_key(int i, int key) {
    if (key < keys[i]) {
      cout << "new key smaller then current key";
      return;
    }
    keys[i] = key;
    while(i > 0 && keys[heap::parent(i)] < keys[i]) {
      swap(&keys[heap::parent(i)], &keys[i]);
      i = heap::parent(i);
    }
}

void max_heap_insert(const int k) {
    heap_size += 1 ;
    keys[heap_size - 1] = INT_MIN;

    heap_increase_key(heap_size - 1, k);
}

\end{lstlisting}

Le due procedure \textbf{non} producono lo stesso \textit{heap}: dato in input $(1,2,3,4,5,6)$, \textit{build\_max\_heap} produce $(6,5,3,4,2,1)$; mentre \ $(6,4,5,1,3,2)$.

\paragraph{Complessità}
Il caso peggiore dell'operazione \textit{build\_max\_heap\_v2} che si verifica quando il vettore è già ordinato in maniera decrescente. Infatti, in questo caso saranno effettuati $n$ \textit{swap}. In questa situazione, il costo $T(n)$ è:

\begin{equation*}
    T(n) = \sum_{i=2}^n \log i = \log (n) + \log (n-1) + \dots + \log (2) = \log (n!)
\end{equation*}

Usando l'approssimazione di \textbf{Stirling}:

\begin{align*}
    T(n) = \log (n!) &\approx \log \left(\sqrt{2\pi n} \left(\frac{n}{e}\right)^n \cdot \left(1 + \Theta\left(\frac{1}{n}\right)\right)\right) \\
    &= \log (\sqrt{2 \pi n}) + \log\left(\frac{n}{e}\right)^n + \log \left(1 + \Theta\left(\frac{1}{n}\right)\right)\\ 
    &= \Theta(n) + \Theta(n \log n) + \Theta\left(\frac{1}{n}\right)
\end{align*}

Usando la proprietà di algoritmi in sequenza(\cref{eq:alg-seq}), si ottiene:

\begin{equation*}
    T(n) = \Theta(n \log n)
\end{equation*}

\clearpage

\subsection{Problema 3}
Sia $A[1\dots n]$ un \textit{array} di $n$ elementi distinti “quasi ordinato”, ovvero in cui ogni elemento dell’\textit{array} si trova entro $k$ slot dalla sua posizione corretta. Definiamo una "inversione": se $i<j$ e $A[i] > A[j]$, allora la coppia $(i,j)$ è detta inversione di $A$. Si implementi un algoritmo che ordina il vettore $A$ in tempo $O(n \log k)$. Suggerimento: si usi un Heap.

\paragraph{Soluzione} Per la soluzione è stato utilizzando un \textit{min heap} (di implementazione analogo al \textit{max heap} mostrato nel problema precedente)  di capacità $k$ e si è sfruttato il fatto che, applicando la procedura di inserimento, l'elemento in posizione $0$ sia il minimo del sotto-array $A[\dots,i, \dots k+i]$. Quindi valida l'ipotesi del vettore "quasi ordinato", l'algoritmo risolutivo è il seguente:


\begin{lstlisting}[caption={Algoritmo per ordinare un vettore "quasi ordinato" mediante heap},captionpos=b]
void sort_k(int arr[], int n, int k) {
  int size = (n == k) ? k : k + 1;

  // si crea un heap
  struct heap h(size);

  // inserimento dei primi k elementi nell'heap
  for (int i = 0; i < size; ++i) {
    h.min_heap_insert(arr[i]);
  }

  int index = 0;
  // estrazione del minimo, rimpiazzo nel vettore iniziale
  // inserimento del prossimo elemento nel vettore heap
  for (int i = k + 1; i < n; i++) {
    arr[index++] = h.extract_min();
    h.min_heap_insert(arr[i]);
  }

  // svuotamento dell'heap con i valori rimasti
  while (!h.empty()) {
    arr[index++] = h.extract_min();
  }
}
\end{lstlisting}


\paragraph{Complessità}
Per il costo $T(n)$ dell'algoritmo, bisogna analizzare singolarmente i $3$ cicli che compongono l'algoritmo.

Il primo ciclo ha un costo pari a:
\begin{equation*}
    T_1(n) = O(k \log k)
\end{equation*}

Il secondo ha un costo pari a:
\begin{equation*}
    T_2(n) = O((n-k)(\log k + \log k)) = O((n-k)\log (k)^2) = O(2(n-k)\log k)
\end{equation*}

Infine, il terzo ha un costo pari a:
\begin{equation*}
    T_3(n) = O(k\log k)
\end{equation*}

Quindi applicando la proprietà in \cref{eq:alg-seq}, il costo totale dell'algoritmo è:

\begin{equation*}
    T(n) = O(\max\{T_1(n), T_2(n), T_3(n)\}) = T_2(n) =  O(2(n-k) \log k) = O(n \log k)
\end{equation*}

\clearpage

\section{Homework}
\subsection{Problema 1}
Sia data una matrice di $N x M$ numeri interi positivi. Si implementi un algoritmo di
\textit{backtracking} per determinare la lunghezza della più lunga sequenza di numeri adiacenti strettamente crescenti (un numero x è adiacente ad y si trova in alto, in alto a destra, in alto a sinistra, a destra, a sinistra, in basso, in basso a destra, in basso a sinistra rispetto ad y).

\paragraph{Soluzione} La soluzione basata su \textit{backtracking} prevede l'esplorazione di tutto lo spazio delle soluzioni. Come mostrato nel listato sottostante, la procedura utilizza una matrice di appoggio in cui l'elemento $A[i][j]=true$ se la cella è parte della soluzione in esame. Questa tecnica è utilizzata per evitare \textit{loop} in quanto previene di ripercorrere due volte la stessa cella nell'arco della stessa soluzione.
All'inizio della funzione, si aggiorna il valore massimo rispetto a quello corrente, si incrementa il valore corrente della soluzione e si marca la cella $(i,j)$ come visitata; infine si cerca la prima direzione ammissibile per trovare una nuova soluzione. Una volta esplorate tutte le soluzioni, la cella viene marcata come non visitata e si decrementa il valore della soluzione corrente.

\begin{lstlisting}[caption={Procedura di calcolo della lunghezza sequenza crescente},captionpos=b]
void _find_increasing_path(int mat[N][M], int n, int m, int i, int j, bool visited[N][M], int& current, int &res) {
  res = max(res, current);

  visited[i][j] = true;
  current ++;

  // sotto
  if (is_valid(mat, n, m, i+1, j, mat[i][j], visited)) {
    _find_increasing_path(mat,n, m, i+1, j, visited, current, res);
  }

  // destra
  if (is_valid(mat, n, m, i, j+1, mat[i][j], visited)) {
    _find_increasing_path(mat,n, m, i, j+1, visited, current, res);
  }

  // sopra
  if (is_valid(mat, n, m, i-1, j, mat[i][j], visited)) {
    _find_increasing_path(mat,n, m, i-1, j, visited, current, res);
  }

  // sinistra
  if (is_valid(mat, n, m, i, j-1, mat[i][j], visited)) {
    _find_increasing_path(mat,n, m, i, j-1, visited, current, res);
  }

  // alto sinistra
  if (is_valid(mat, n, m, i-1, j-1, mat[i][j], visited)) {
    _find_increasing_path(mat,n, m, i-1, j-1, visited, current, res);
  }

  // alto destra
  if (is_valid(mat, n, m, i-1, j+1, mat[i][j], visited)) {
    _find_increasing_path(mat,n, m, i-1, j+1, visited, current, res);
  }
  
  // basso sinistra
  if (is_valid(mat, n, m, i+1, j-1, mat[i][j], visited)) {
    _find_increasing_path(mat,n, m, i+1, j-1, visited, current, res);
  }
  
  // basso destra
  if (is_valid(mat, n, m, i+1, j+1, mat[i][j], visited)) {
    _find_increasing_path(mat,n, m, i+1, j+1, visited, current, res);
  }
  
  visited[i][j] = false;
  current --;
}

\end{lstlisting}
Questa procedura ricorsiva viene richiamata da un \textit{wrapper} che si limita ad inizializzare le variabili. 

\paragraph{Complessità}
La complessità è data dalla ricorrenza:
\begin{equation*}
    T(n) = 8T(n-1) + 1
\end{equation*}

Si osserva che la $T(n)$ è una ricorrenza lineare di ordine costante con $a=8$ e $\beta = 0$. Pertanto, è possibile affermare che la complessità dell'algoritmo è:
\begin{equation*}
    T(n) = \Theta(a^n n^\beta) = \Theta(8^n)
\end{equation*}

\clearpage
\subsection{Problema 2}

Sia data un vettore di n numeri interi positivi. Si considerino le sotto-sequenze i cui valori sono ordinati in senso crescente (si assume che non ci siano duplicati). Si scriva un
programma per trovare la somma massima tra le somme degli elementi di tali sotto-
sequenza. Ad esempio, se l'array di input è $\{3, 4, 5, 10\}$, l'output dovrebbe essere $22 = 3 + 4 + 5 + 10$; se l'array di input è $\{10, 5, 4, 3\}$, l'output dovrebbe essere $10$.

\paragraph{Soluzione} Il problema è stato risolto attraverso la programmazione dinamica. In particolare, dovendo calcolare solo la somma, basta utilizzare il vettore $dp$ a una dimensione. L'algoritmo risolutivo è mostrato di seguito:

\begin{lstlisting}[caption={Calcolo della sottosequenza crescente di somma massima},captionpos=b]

int max_sum_increasing_subsequence(const int arr[], const int n) {

    int res = INT_MIN;
    int dp[n];

    // inizializzazione vettore dp
    for (size_t i = 0; i < n; ++i) {
        dp[i] = arr[i];
    }

    // calcolo della somma massima per la sottosequenza che termina in i-esima posizione
	for (size_t i = 1; i < n; ++i) {
    for (size_t j = 0; j < i; ++j) {
      // sono in una sottosequenza crescente?
      if (arr[i] > arr[j]) {
        // incremento la somma
        dp[i] = max(dp[i], dp[j] + arr[i]);
      }
    }
  }

  // calcolo del valore massimo
	for (size_t i = 0; i < n; ++i) {
		if (res < dp[i]) {
		    res = dp[i];
    }
  }
	return res;
}
\end{lstlisting}
Il vettore $dp$ contiene alla posizione $i-esima$ il valore della somma della sotto-sequenza che termina nell'elemento $i-esimo$ dell'array di input.

\paragraph{Complessità} 
L'algoritmo ha un costo $T(n)$ dato dalla somma dei 3 cicli che lo compongono (\cref{eq:alg-seq}). Il primo ciclo inizializza il vettore $dp$ con l'elemento omologo del vettore e ha un costo di: 
\begin{equation*}
    T_1(n) = O(n)
\end{equation*}

Il secondo blocco dell'algoritmo ha due cicli innestati che, sfruttando la proprietà in \cref{eq:nested-loops}, ha una complessità:

\begin{equation*}
    T_2(n) = O(n^2)
\end{equation*}

Il terzo ciclo ha un costo totale di:
\begin{equation*}
    T_3(n) = O(n)
\end{equation*}

Quindi il costo totale è:
\begin{equation*}
    T(n) = \max\{T_1(n), T_2(n), T_3(n)\} = T_2(n) = O(n^2)
\end{equation*}

Infine, la complessità spaziale è $S(n) = O(n)$ data dall'utilizzo del vettore $dp$.
\clearpage

\subsection{Problema 3}
Un numero può sempre essere rappresentato come somma di quadrati di altri numeri.
Infatti, al limite, il valore 1 è un quadrato e possiamo sempre rappresentare un qualsiasi
numero come ($1\cdot 1 + 1\cdot 1 + 1\cdot 1 + \dots$).
Dato un numero n, si implementi un algoritmo che trova il numero minimo di quadrati la cui
somma è $n$. Ad esempio, se $n = 100$, questo valore si può ottenere sommando $5^2 + 5^2 + 5^2 + 5^2$ , quindi usando $4$ quadrati, ma anche come $10^2$, quindi usando un solo quadrato.

\paragraph{Soluzione}
La soluzione del problema prevede l'utilizzo della programmazione dinamica. In particolare, ci si limita a costruire la soluzione a partire dalla relazione ricorsiva:

\begin{equation*}
    DP[i]=\begin{cases}
        \min(DP[i], 1 + DP(i-k^2))& i > 1 \\
        1 & i \leq 1 \\
    \end{cases}
\end{equation*}
con $1\leq k \leq i \leq n$. L'approccio utilizzato è di tipo \textit{bottom-up}. Per ogni numero $i \leq n$, l'algoritmo lo può esprimere come somma di quadrati sottraendo un numero  $j\leq i$ al quadrato $(i - j^2)$ e quindi riconducendosi a un sotto-problema precedente.

\begin{lstlisting}[caption={Calcolo del numero minimi di quadrati per rappresentare un numero $n$},captionpos=b]
int min_square(const int n) {
  
  // inizializzo il vettore dp
  int dp[n+1];

  dp[0] = 0;
  dp[1] = 1;
  for(int i = 2; i <= n; ++i) {
    dp[i] = INT_MAX;
    for(int j = 1; (j*j) <= i; ++j) {
      // applico la formula ricorsiva
      dp[i] = min(dp[i], 1 + dp[i - (j*j)]);
    }
  }
  return dp[n];
}
\end{lstlisting}

\paragraph{Complessità}
Innanzitutto, la complessità spaziale dell'algoritmo è $\Theta(n)$ in quanto viene creato un vettore con $n+1$ elementi. Mentre, il costo computazionale è $O(n\sqrt{n})$, ottenuto applicando la proprietà dei cicli innestati (\cref{eq:nested-loops}): il primo ciclo ha una complessità $O(n)$ mentre il secondo $O(\sqrt{n})$.

\clearpage

\section{Homework}
\subsection{Problema}
La S.S.C. Napoli è interessata a massimizzare i guadagni derivanti dagli ingressi allo stadio.
Ti hanno assunto come consulente per aumentare le vendite dei “biglietti per gruppi” di
tifosi. Bisogna risolvere il seguente problema: quando un gruppo vuole vedere una partita,
tutti i membri del gruppo devono avere un posto, altrimenti se ne vanno. Quindi, poiché non
vi possono essere gruppi “parziali”, le gradinate spesso non sono piene. C'è ancora spazio
disponibile, ma non abbastanza per un intero gruppo. In tal caso, il gruppo non può sedersi,
e il Napoli perde soldi. Si sviluppi un algoritmo per decidere quali gruppi far entrare per
massimizzare le vendite.

\paragraph{Soluzione} L'esercizio è un problema di ottimizzazione in cui si vuole massimizzare il profitto ottenuto facendo entrare gruppi più grandi di persone. Infatti, può essere ricondotto al \textit{knapsack problem} noto in letteratura e risolvibile all'ottimo attraverso la programmazione dinamica. In particolare, il problema può essere modellato come segue:
\begin{itemize}
    \item $C$ capacità dello stadio;
    \item $N$ gruppi di tifosi. L'$i-esimo$ gruppo è costituito da $w[i]$ persone;
    \item Ad ogni gruppo $i$ di tifosi è associato un valore di profitto $p[i]$ dato dall'incasso della vendita dei biglietti. Per semplicità, si assume che il costo del biglietto per tifoso sia unitario e, quindi $p[i]$ risulta pari al numero di persone nel gruppo $i-esimo$.
\end{itemize}

Il problema dello zaino può essere formulato come segue:
\begin{equation*}
    DP[i][c]=\begin{cases}
        \max(DP[i-1][c], DP[i-1][c-w[i]] + p[i])& i > 1 \; \land\\
        &c - w[i] \geq 0 \\
        0 & i = 0 \lor c = 0 \\
    \end{cases}
\end{equation*}
con $0 \leq c \leq C$ e $0 \leq i \leq N$. Il caso base è dovuto dal fatto che se non ci sono gruppi di tifosi ($i = 0$) oppure lo stadio non ha posti a sedere ($c = 0$), l'incasso è pari a $0$.


\begin{lstlisting}[caption={Implementazione algoritmo dello zaino con ricostruzione della soluzione},captionpos=b]
vector<int> knapsack(const int w[], const int p[], int n, int c) {
  int dp[n+1][c+1];
  int res = 0;
  vector<int> sol;
  // capienza = 0
  for (size_t i = 0; i <= n; ++i) {
    dp[i][0] = 0;
  }
  // tifosi = 0
  for (size_t i = 0; i <= c; ++i) {
    dp[0][i] = 0;
  }

  for(size_t i = 1; i <= n; ++i){
    for (size_t j = 1; j<= c; ++j) {
      // se ho capienza metto l'oggetto nello zaino
      if (w[i-1] <= j) {
        dp[i][j] = max(dp[i-1][j], dp[i-1][j-w[i-1]] + p[i-1]);
      } else {
        // altrimenti il profitto rimane invariato
        dp[i][j] = dp[i-1][j];
      }
    }
  }

  res = dp[n][c];
  for (size_t i = n, j = c; i > 0 && res > 0; i--) {
    // se il risultato si trova in dp[i-1][j]
    // l'elemento non non fa parte della soluzione
    if (res == dp[i-1][j]) {
     continue;   
    } 
    // altrimenti si salva il risultato nell'array
    sol.push_back(i-1);
             
    // si decrementa res e j per saltare ad un altra colonna
    res -= p[i-1];
    j -= w[i-1];
  }

  reverse(sol.begin(), sol.end());
  return sol;
}

\end{lstlisting}

La funzione $knapsack$ accetta in input il vettore dei pesi, quello dei profitti, la capacità dello stadio e la dimensione del problema e restituisce la soluzione formata da un vettore contente gli indici dei gruppi da far entrare nello stadio per massimizzare i profitti.

\paragraph{Complessità}
L'algoritmo proposto ha complessità spaziale $S(n) = \Theta(n\cdot C)$ dato dal costo di memorizzazione della tabella $dp$. Analogamente, la complessità temporale è $T(n) = \Theta(n\cdot C)$ data dai cicli innestati per lavorare sulla matrice.

\end{document}
